\chapter{Recorridos eulerianos y hamiltonianos}
\label{md:recorridos}

\section{Caminos y conectividad}
\label{md:caminos}

Un \textbf{camino} es una sucesión de vértices y aristas incidentes
que permite pasar de cada vértice al siguiente.
Su \textbf{longitud} es el número de aristas recorridas; se permiten
repeticiones. En multigrafos hay que especificar qué arista se utiliza.

Un camino es \textbf{simple} si no repite vértices y \textbf{cerrado}
si sus extremos coinciden. Un \textbf{circuito} es un camino cerrado
de longitud positiva sin repetición de aristas. Un \textbf{ciclo}
es un circuito sin repetición de vértices salvo la coincidencia
del primero con el último. En un grafo simple un ciclo tiene longitud
al menos tres; dos aristas paralelas pueden formar uno de longitud dos.

Un grafo es \textbf{conexo} si cualquier pareja de vértices está unida
por un camino. Si hay un camino entre vértices distintos, se puede
obtener uno simple eliminando los tramos cerrados producidos por
vértices repetidos.
La relación de estar unidos por un camino, incluyendo cada vértice
consigo mismo, es de equivalencia. Sus clases forman las
\textbf{componentes conexas}; no hay aristas entre componentes distintas.

\section{Euler: recorrer todas las aristas}
\label{md:euler-grafos}

Un \textbf{camino euleriano} recorre cada arista exactamente una vez.
Si es cerrado, es un \textbf{circuito euleriano}.
Para un grafo no dirigido conexo con al menos una arista:
\begin{itemize}
 \item Existe un circuito euleriano si y solo si todos los grados son pares.
 \item Existe un camino euleriano abierto si y solo si hay exactamente
       dos vértices de grado impar; estos son sus extremos.
\end{itemize}
La condición de conectividad impide que queden aristas en otra componente.
La necesidad de la paridad se debe a que, en un vértice intermedio,
cada entrada se empareja con una salida; los extremos de un recorrido
abierto tienen una incidencia sin pareja.

Para \textbf{construir un circuito}, se parte de un vértice y se
siguen aristas aún no usadas hasta volver al inicio.
Si quedan aristas, la conectividad permite encontrar en el circuito
un vértice incidente con alguna de ellas. Desde él se construye otro
circuito usando las restantes y se inserta en el anterior.
Se repite hasta agotar todas las aristas.
Con dos grados impares se puede añadir una arista entre ellos,
obtener un circuito y suprimir esa arista para abrirlo.

\section{Hamilton: visitar todos los vértices}
\label{md:hamilton}

Un \textbf{camino hamiltoniano} es simple y visita todos los vértices.
Un \textbf{ciclo hamiltoniano} pasa por todos ellos y vuelve al inicio,
sin otras repeticiones. Un grafo con tal ciclo es hamiltoniano;
para grafos simples se considera $|V|\ge3$.
El objetivo se refiere a los vértices; no requiere usar todas las aristas.

Si un grafo tiene ciclo hamiltoniano y $\varnothing\ne U\subseteq V$,
al borrar los vértices de $U$ el número de componentes restantes
no puede superar $|U|$.
Un ciclo queda dividido en a lo sumo $|U|$ tramos; otras aristas
pueden unirlos, pero no crear más componentes.
Es una condición necesaria, no suficiente.

El \textbf{teorema de Dirac} da una condición suficiente:
un grafo simple con $n\ge3$ vértices y
\[
 \gr(v)\ge\frac n2\quad\text{para todo }v
\]
tiene ciclo hamiltoniano.

Para \emph{caminos} hamiltonianos se dispone de otro criterio:
si un grafo simple con $n>1$ cumple
\[
 \gr(u)+\gr(v)\ge n-1\quad\text{para todo par }u\ne v,
\]
entonces tiene un camino hamiltoniano. En particular, basta
$\gr(v)\ge(n-1)/2$ para todos los vértices.
Estos criterios son suficientes: incumplirlos no descarta por sí solo
un recorrido hamiltoniano.
